\section{Results}

This section evaluates the proposed bird tracking pipeline from three aspects:
tracking metrics, comparison with a baseline method, and computational efficiency.
All experiments are conducted on the validation set with ground-truth annotations
(\texttt{val\_mot}).

\subsection{Metric}
\label{subsec:metric}

We evaluate tracking performance using standard MOT metrics, including
\textbf{MOTA}, \textbf{IDF1}, \textbf{ID Switches (IDs)}, \textbf{False Positives (FP)},
and \textbf{False Negatives (FN)}, computed with the \texttt{motmetrics} toolkit
using an IoU threshold of 0.5.

Qualitative results show that the proposed pipeline is able to recover meaningful
bird trajectories in many sequences.
Tracks are often temporally consistent, especially in scenes with moderate motion
and relatively stable backgrounds.
This supports \textit{Hypothesis 1}, demonstrating the feasibility of a classical
computer vision pipeline for bird tracking.

However, quantitative performance is significantly affected by a large number of
false positives.
Background regions with strong motion are sometimes propagated as valid tracks,
which leads to degraded MOTA and IDF1 scores.
This highlights that false positive suppression is a key limitation of the current system.

\subsection{Comparison}
\label{subsec:comparison}

We compare our method with NetTrack, a deep-learning-based tracking framework.
While NetTrack generally achieves higher MOTA and IDF1 scores, the proposed method
produces visually competitive tracking results on several sequences.
Bird trajectories are clearly captured and maintained when motion is smooth,
and identity switches remain limited in short- to mid-term tracking.

The performance gap is mainly caused by the stronger detection robustness of NetTrack,
which suppresses background clutter more effectively.
These results support \textit{Hypothesis 2}: deep fine-grained tracking performs better
in highly dynamic and cluttered scenes, while classical methods can still provide
reasonable tracking behavior under simpler conditions.

\subsection{Efficiency}
\label{subsec:efficiency}

The proposed pipeline is fully based on classical computer vision techniques,
including motion compensation, Lucas--Kanade optical flow, and greedy IoU-based
data association.
It runs entirely on CPU, requires no training data, and maintains low memory usage.

In practice, the proposed method is significantly faster than NetTrack,
especially on long sequences.
By avoiding deep neural network inference, it achieves much lower runtime while
sacrificing tracking accuracy.
This confirms \textit{Hypothesis 3}, highlighting a clear trade-off between efficiency
and performance.